\chapter{Dispersion et corrélation}\label{ch:s2:dispersion-and-correlation}

Une option sur indice est un pari sur la volatilité de l'indice, qui dépend de la volatilité de ses composantes et de la façon dont elles évoluent ensemble. Cboe publie la
corrélation que les prix d'options impliquent parmi les 50 plus grandes valeurs de l'S\&P 500~: de 2006 à 2026, elle a valu en moyenne 37, a atteint 96.6 le 24 octobre 2008, puis est tombée
à 2.93 en juillet 2024. Driessen, Maenhout et Vilkov ont établi que le risque de corrélation est tarifé~: le vendre, en vendant la volatilité de l'indice contre
la volatilité des composantes, rapportait un alpha élevé avant frictions et aucun après celles-ci. Sur le marché synthétique de ce chapitre, la même transaction affiche un ratio de Sharpe de 0.65
rendue neutre en vega, et un swap de corrélation vendeur pur 1.74, jusqu'à ce qu'un krach fasse passer la corrélation de 0.12 à 0.83 en un mois. La construction est
\texttt{firm.dispersion}.

\section{Volatilité de l'indice à partir des volatilités des composantes}

La variance d'un indice est la somme pondérée des variances et des covariances de ses composantes~: $\sigma_I^2 = \sum_i w_i^2\sigma_i^2 + \sum_{i \ne j} w_i w_j
\rho_{ij}\sigma_i\sigma_j$. En remplaçant chaque corrélation par paires par une moyenne unique $\bar\rho$ et en résolvant, on obtient la corrélation moyenne qui rend les
volatilités des composantes cohérentes avec celle de l'indice~:
\[
\bar\rho = \frac{\sigma_I^2 - \sum_i w_i^2\sigma_i^2}{\sum_{i \ne j} w_i w_j \sigma_i\sigma_j}.
\]
Les indices de corrélation implicite de Cboe appliquent cette décomposition aux volatilités implicites du S\&P 500 et de ses 50 plus grandes composantes.

\begin{definition}[Corrélation réalisée]\label{def:s2:dispersion-and-correlation:realised}
La \emph{corrélation réalisée}\index{corrélation réalisée} est la corrélation moyenne entre les composantes d'un indice, calculée à partir des variances réalisées sur une
période, avec la même décomposition que la corrélation implicite~: la variance réalisée de l'indice comparée aux variances réalisées des composantes.
\end{definition}

\texttt{firm.dispersion} (le \cref{lst:s2:dispersion-and-correlation:legs}) calcule les deux. Sur le marché synthétique, l'indice est la moyenne à pondération égale de
trente composantes~; la variance implicite des composantes est leur variance attendue multipliée par 1.10, et celle de l'indice par 1.30, si bien que les options sur l'indice sont plus chères
que celles de ses composantes, et la corrélation implicite se situe au-dessus de ce qui se réalise~: 0.311 en moyenne contre une \omterm{def:s2:dispersion-and-correlation:realised}{corrélation réalisée} de 0.252, et plus élevée dans 77~\% des mois.

\section{La prime de risque de corrélation}

\begin{definition}[Prime de risque de corrélation]\label{def:s2:dispersion-and-correlation:premium}
La \emph{prime de risque de corrélation}\index{prime de risque de corrélation} est l'écart entre la corrélation moyenne impliquée par les prix d'options sur l'indice et sur ses
composantes et la \omterm{def:s2:dispersion-and-correlation:realised}{corrélation réalisée} ensuite~; elle est positive en moyenne, car la corrélation monte pendant les krachs, au moment où la diversification est la plus nécessaire.
\end{definition}

Driessen, Maenhout et Vilkov ont utilisé les options sur l'S\&P 100 et sur toutes ses composantes~: le risque de corrélation était tarifé, il expliquait la
coupe transversale des rendements des options sur indice et sur valeurs isolées, et une stratégie qui le vendait avait un alpha élevé pour un investisseur sans frictions. Avec des coûts de transaction réalistes, la
prime ne pouvait pas être exploitée, ce qu'ils ont interprété comme une limite à l'arbitrage. La comparaison de Cboe elle-même avec une \omterm{def:s2:dispersion-and-correlation:realised}{corrélation réalisée} lissée
des 50 premières valeurs conclut que la \omterm{def:s2:dispersion-and-correlation:realised}{corrélation réalisée} est en général un peu inférieure à la corrélation implicite.

L'indice réel (la \cref{fig:s2:dispersion-and-correlation:cor}) montre pourquoi vendre la corrélation est inconfortable~: le COR1M a valu en moyenne 37.0 de janvier 2006 à
septembre 2026, avec des quantiles à 5~\% et 95~\% de 10.8 et 69.7, et il a valu en moyenne 70.4 en octobre 2008 et 72.5 en mars 2020. Le 22 septembre 2026, il se situait à
7.59~; son plus bas, 2.93, a été atteint le 12 juillet 2024. L'indice de dispersion de Cboe sur l'S\&P 500, publié depuis
juin 2014, a valu en moyenne 25.5 et culminé à 58.9 le 18 mars 2020.

\begin{omfigure}
\begin{tikzpicture}
\begin{axis}[width=11cm,height=5.2cm,xlabel={\small année},ylabel={\small COR1M, moyenne mensuelle},tick label style={font=\footnotesize},
  xmin=2006,xmax=2027,ymin=0,ymax=100,grid=major,grid style={gray!25},table/col sep=comma,/pgf/number format/1000 sep={}]
\addplot[omThm,thick] table[x=year,y=cor1m]{figdata/strategies-2/02-dispersion-and-correlation/cor1m.csv};
\end{axis}
\end{tikzpicture}

\omcaption{Indice de corrélation implicite à un mois de Cboe (la corrélation moyenne entre les 50 plus grandes valeurs de l'S\&P 500, implicite dans les prix des options sur l'indice et sur les valeurs
isolées), moyennes mensuelles, de janvier 2006 à septembre 2026. Dérivé de l'historique d'indices de Cboe~; la série brute n'est pas redistribuée. Données~:
\texttt{s2\_fetch\_cor}.}\label{fig:s2:dispersion-and-correlation:cor}
\end{omfigure}

\section{Trades de dispersion et leurs pondérations}

Un trade de dispersion (livre~5, chapitre~17) vend la volatilité de l'indice et achète celle des composantes. Les portefeuilles du chapitre utilisent des swaps de variance à un mois, reconstruits
chaque mois (le \cref{lst:s2:dispersion-and-correlation:books})~: vendeur d'une unité de variance de l'indice, acheteur de la variance des composantes, sur des notionnels fixés de l'une de deux façons.

\begin{definition}[Pondération de la dispersion]\label{def:s2:dispersion-and-correlation:weighting}
La \emph{pondération de la dispersion}\index{pondération de la dispersion} est le choix des notionnels des composantes dans un trade de dispersion~: égaux à leurs poids dans l'indice, de sorte
que le trade soit vendeur de corrélation et acheteur de la variance propre aux composantes~; ou ajustés pour que le vega total des composantes égale celui de l'indice
(neutre en vega), de sorte qu'une hausse commune de la volatilité laisse le portefeuille à plat et que son P\&L provienne principalement de la corrélation.
\end{definition}

\begin{center}
\small
\begin{tabular}{@{}lcccc@{}}
\toprule
239 mois & ratio de Sharpe & asymétrie & pire mois (écart-type) & mois haussiers\\
\midrule
composantes aux poids de l'indice & $-0.02$ & 2.67 & $-1.6$ & 36~\%\\
neutre en vega & 0.65 & $-5.87$ & $-11.3$ & 70~\%\\
swap de corrélation vendeur & 1.74 & $-1.67$ & $-6.0$ & 77~\%\\
\bottomrule
\end{tabular}
\end{center}

Avec une pondération selon les poids de l'indice, le portefeuille est plus long en variance des composantes qu'il n'est court en variance de l'indice, et il paie la prime propre aux composantes~: il perd un peu en
moyenne et gagne pendant les krachs (asymétrie positive). Rendu neutre en vega, il rapporte la prime de corrélation avec la forme habituelle des \omterm{def:s2:harvesting-the-variance-risk-premium:short}{stratégies vendeuses de volatilité}. Le
swap de corrélation, qui verse directement la corrélation implicite moins la \omterm{def:s2:dispersion-and-correlation:realised}{corrélation réalisée}, isole le mieux la prime~; son ratio de Sharpe est le plus élevé, et son pire mois
reste encore à six écarts-types.

\section{Quand la corrélation tend vers un}

Les trois krachs introduits dans les données sont les tests (la \cref{fig:s2:dispersion-and-correlation:books}). Dans le premier, la corrélation implicite valait 0.12 au début du mois
et la \omterm{def:s2:dispersion-and-correlation:realised}{corrélation réalisée} 0.83~: le portefeuille neutre en vega a perdu 11.3 de ses écarts-types mensuels, et le swap de corrélation 0.71 d'une unité. Dans le deuxième,
0.12 contre 0.57 (3.0 écarts-types et 0.45). Dans le troisième, qui a commencé alors que la corrélation implicite valait déjà 0.73, la perte est restée inférieure à
un écart-type. Un krach est un événement de corrélation~: toutes les composantes chutent ensemble, l'indice baisse autant que ses composantes, et la
diversification que l'option sur indice était tarifée pour ne pas offrir arrive alors.

\begin{omfigure}
\begin{tikzpicture}
\begin{axis}[width=11cm,height=5.2cm,xlabel={\small année},ylabel={\small P\&L cumulé (unités d'écart-type)},tick label style={font=\footnotesize},
  xmin=0,xmax=20,ymin=-20,ymax=130,grid=major,grid style={gray!25},legend pos=north west,legend style={font=\scriptsize},table/col sep=comma]
\addplot[omDef,thick] table[x=year,y=vega]{figdata/strategies-2/02-dispersion-and-correlation/books.csv};
\addplot[omThm,thick] table[x=year,y=corrswap]{figdata/strategies-2/02-dispersion-and-correlation/books.csv};
\legend{dispersion neutre en vega,swap de corrélation vendeur}
\end{axis}
\end{tikzpicture}

\omcaption{Le portefeuille de dispersion neutre en vega et le swap de corrélation vendeur synthétiques, chacun cumulé en unités de son écart-type mensuel, sur
vingt ans, avec des krachs aux années 6.0, 12.7 et 17.5. Données~: \texttt{s2\_dispersion.monthly}.}\label{fig:s2:dispersion-and-correlation:books}
\end{omfigure}

\section{Fiches de stratégie}

\begin{strategyfile}{Dispersion pondérée par le vega}\label{strat:s2:dispersion-and-correlation:vega}
\sfield{Qui paie, et pourquoi} Acheteurs de protection sur indice, qui paient pour la corrélation que font surgir les krachs.
\sfield{Instruments et plateformes} Options sur indice ou swaps de variance, contre des options sur titres individuels ou des swaps de variance sur les plus gros membres.
\sfield{Signal} Corrélation implicite contre corrélation prévisionnelle.
\sfield{Dimensionnement et exécution} Vega des membres apparié à celui de l'indice~; roll mensuels ou trimestriels~; membres limités aux titres liquides.
\sfield{Coûts} Fourchettes des options sur titres individuels, principale friction.
\sfield{Comment elle disparaît} Krachs~; frictions qui absorbent la prime.
\sfield{Horizon, capacité, infrastructure} Quelques mois~; surfaces de volatilité sur titres individuels pour plusieurs dizaines de noms.
\sfield{Backtest honnête} Coûts des options sur titres individuels~; mois de krach dans l'échantillon.
\sfield{Sources} Driessen, Maenhout et Vilkov (2009)~; ce chapitre~: 0.65 et un mois de krach à 11 écarts-types.
\end{strategyfile}

\begin{strategyfile}{Dispersion pondérée par le theta}\label{strat:s2:dispersion-and-correlation:theta}
\sfield{Qui paie, et pourquoi} Comme pour la dispersion pondérée par le vega.
\sfield{Instruments et plateformes} Straddles sur indice et sur titres individuels.
\sfield{Signal} Comme précédemment.
\sfield{Dimensionnement et exécution} Theta des membres apparié à celui de l'indice, de sorte que l'érosion temporelle du portefeuille soit neutre et que son P\&L soit un gamma contre gamma.
\sfield{Coûts} Couverture en delta de nombreuses options.
\sfield{Comment elle disparaît} Comme précédemment.
\sfield{Horizon, capacité, infrastructure} De quelques semaines à quelques mois~; un moteur de couverture.
\sfield{Backtest honnête} Coûts de couverture pour chaque membre.
\sfield{Sources} Aucun chiffre de performance vérifié.
\end{strategyfile}

\begin{strategyfile}{Swap de corrélation vendeur}\label{strat:s2:dispersion-and-correlation:corrswap}
\sfield{Qui paie, et pourquoi} Acheteurs de protection contre la corrélation, souvent des desks de produits structurés (chapitre~28).
\sfield{Instruments et plateformes} Swaps de corrélation de gré à gré.
\sfield{Signal} Strike comparé à la corrélation prévisionnelle.
\sfield{Dimensionnement et exécution} Dimensionné sur la perte du mois de krach.
\sfield{Coûts} Fourchettes des dealers~; exposition de contrepartie.
\sfield{Comment elle disparaît} Krachs.
\sfield{Horizon, capacité, infrastructure} De quelques mois à un an~; une relation de gré à gré.
\sfield{Backtest honnête} \omterm{def:s2:dispersion-and-correlation:realised}{Corrélation réalisée} mesurée comme le contrat la définit.
\sfield{Sources} Ce chapitre~: 1.74 avant coûts, une perte de 0.71 unité lors du premier krach.
\end{strategyfile}

\begin{strategyfile}{Dispersion sectorielle}\label{strat:s2:dispersion-and-correlation:sector}
\sfield{Qui paie, et pourquoi} Comme pour la dispersion sur indice, au sein d'un secteur.
\sfield{Instruments et plateformes} Options sur ETF sectoriels contre les options de leurs plus gros membres.
\sfield{Signal} Corrélation implicite sectorielle comparée à son historique.
\sfield{Dimensionnement et exécution} Neutre en vega au sein de chaque secteur.
\sfield{Coûts} Fourchettes plus larges sur les options sectorielles.
\sfield{Comment elle disparaît} Chocs touchant l'ensemble du secteur.
\sfield{Horizon, capacité, infrastructure} Mois~; capacité réduite.
\sfield{Backtest honnête} Liquidité des options sectorielles.
\sfield{Sources} Corrélations implicites sectorielles de Cboe (comme source de données)~; aucun chiffre de performance vérifié.
\end{strategyfile}

\section{Tutoriel~: quand tout bouge ensemble}\label{tut:s2:dispersion-and-correlation:tut}

\begin{tutorial}
\textbf{Objectif.} Mesurer la corrélation implicite et réalisée sur l'indice synthétique et sur ses membres, faire passer trois portefeuilles de dispersion à travers les
krachs introduits, et lire les indices de corrélation et de dispersion de Cboe. \textbf{État final~:} le tableau et les deux figures.
\begin{tutsteps}
\item \textbf{Corrélation et jambes}.
\omcode{code/firm/dispersion/firm_dispersion.py}{21}{47}{Corrélation moyenne, variance réalisée et P\&L de la dispersion.}\label{lst:s2:dispersion-and-correlation:legs}
\item \textbf{Les portefeuilles mensuels}.
\omcode{code/strategies-2/02-dispersion-and-correlation/python/s2_dispersion.py}{35}{58}{Corrélations implicite et réalisée et trois portefeuilles, mois par mois.}\label{lst:s2:dispersion-and-correlation:books}
\item \textbf{Exécuter} \texttt{s2\_fetch\_cor.py} une fois, puis \texttt{monthly()}, \texttt{crashes()} et \texttt{fig\_dispersion.py}.
\end{tutsteps}
\textbf{Pour aller plus loin.} Ne négocier que les vingt plus gros membres et mesurer le suivi de l'indice~; ajouter un skew de corrélation afin que la corrélation de
krach soit valorisée~; couvrir le portefeuille neutre en vega avec une petite position longue en corrélation avant que les périodes calmes ne prennent fin.
\end{tutorial}

\section{Construction~: dispersion}\label{bld:s2:dispersion-and-correlation:dispersion}

\begin{build}
\bfield{Rôle} Corrélation moyenne implicite et réalisée, portefeuilles de dispersion sur swaps de variance et swaps de corrélation.
\bfield{Interface} \texttt{average\_corr(var\_index, var\_members, w)}, \texttt{realised\_var(R, start, n)},
\texttt{dispersion\_pnl(iv\_index, iv\_members, rv\_index, rv\_members, w, vega\_neutral)}, \texttt{corr\_swap\_pnl(implied, realised)}.
\bfield{Règles} La même décomposition pour la corrélation implicite et réalisée~; les jambes des swaps de variance en unités de variance annuelles.
\bfield{Tests d'acceptation} \texttt{code/firm/dispersion/tests/}~: une corrélation commune retrouvée exactement~; P\&L nul lorsque chaque jambe réalise sa
variance implicite~; les notionnels neutres en vega, calculés à la main.
\bfield{Pour aller plus loin} Jambes de straddles avec couverture~; skew de corrélation~; paniers sectoriels.
\end{build}

\begin{omsources}
\item J.~Driessen, P.~J.~Maenhout et G.~Vilkov, «~The price of correlation risk: evidence from equity options~», \emph{Journal of Finance} 64(3), 2009.
\item Cboe, méthodologie de l'indice S\&P 500 Implied Correlation Index et notes sur la corrélation implicite.
\item Cboe Global Markets, historiques quotidiens de COR1M et DSPX.
\end{omsources}

\section{Exercices}

\begin{exercise}[$\star$]\label{exo:s2:dispersion-and-correlation:1}
Un indice de nombreux membres à pondération égale a une volatilité implicite de 20~\% et chaque membre une volatilité de 30~\%. Calculer la corrélation moyenne que cela implique.
\end{exercise}

\begin{exercise}[$\star$]\label{exo:s2:dispersion-and-correlation:2}
Pour un trade de dispersion neutre en vega, avec une volatilité implicite de l'indice de 20~\% et des membres à 30~\%, calculer le facteur par lequel les notionnels des membres
sont ajustés à partir de leurs poids dans l'indice.
\end{exercise}

\begin{exercise}[$\star$]\label{exo:s2:dispersion-and-correlation:3}
Expliquer pourquoi le portefeuille, construit sur les poids de l'indice, présente une asymétrie positive.
\end{exercise}

\begin{exercise}[$\star\star$]\label{exo:s2:dispersion-and-correlation:4}
Un swap de corrélation vendeur a été fixé à un strike de 0.12 alors que la \omterm{def:s2:dispersion-and-correlation:realised}{corrélation réalisée} vaut 0.83. Calculer la perte par unité et expliquer pourquoi le strike semblait raisonnable au moment de sa fixation.
\end{exercise}

\begin{exercise}[$\star\star$]\label{exo:s2:dispersion-and-correlation:5}
Expliquer pourquoi Driessen et ses coauteurs ont pu mettre en évidence une \omterm{def:s2:dispersion-and-correlation:premium}{prime de risque de corrélation} valorisée qui ne pouvait pas être exploitée.
\end{exercise}

\begin{exercise}[$\star\star$]\label{exo:s2:dispersion-and-correlation:6}
Le COR1M est tombé à 2.93 en juillet 2024. Expliquer ce qu'une faible corrélation implicite révèle sur les options sur indice
par rapport aux options sur titres individuels, puis préciser le risque qu'implique la vente de corrélation dans ce contexte.
\end{exercise}

\begin{exercise}[$\star\star\star$]\label{exo:s2:dispersion-and-correlation:7}
\emph{Programmation.} Exécuter \texttt{monthly(0.30)}, qui fixe la prime des membres égale à celle de l'indice. Que se passe-t-il pour la prime de corrélation et
pour les books~?
\end{exercise}

\begin{exercise}[$\star\star\star$]\label{exo:s2:dispersion-and-correlation:8}
\emph{Trouver la faille.} «~La corrélation implicite a dépassé la \omterm{def:s2:dispersion-and-correlation:realised}{corrélation réalisée} dans 77~\% des mois, donc un swap de corrélation vendeur est une transaction à forte probabilité que l'on peut
faire levier.~»
\end{exercise}

\section{Problème~: quand tout bouge ensemble}

\begin{problem}[{Problème du week-end --- vendre la corrélation}]\label{pb:s2:dispersion-and-correlation:1}
L'indice synthétique du chapitre, les indices de Cboe et les données publiques.

\medskip\noindent\textbf{Partie I --- Corrélation.}
\begin{enumerate}
\item Écrire la formule de la corrélation moyenne et l'expliquer.
\item Définir la \omterm{def:s2:dispersion-and-correlation:realised}{corrélation réalisée} et la \omterm{def:s2:dispersion-and-correlation:premium}{prime de risque de corrélation}.
\item Comment Cboe calcule-t-il son indice de corrélation implicite~?
\item Que Driessen, Maenhout et Vilkov ont-ils mis en évidence~?
\end{enumerate}

\medskip\noindent\textbf{Partie II --- Les indices réels.}
\begin{enumerate}[resume]
\item Donner la moyenne et l'amplitude du COR1M, ainsi que ses niveaux de 2008 et de 2020.
\item Que mesure le DSPX, et quelles étaient ses statistiques~?
\item Pourquoi un plus bas historique de la corrélation implicite n'est-il pas une opération gratuite~?
\item Qu'a constaté Cboe en comparant la corrélation implicite et la \omterm{def:s2:dispersion-and-correlation:realised}{corrélation réalisée}~?
\end{enumerate}

\medskip\noindent\textbf{Partie III --- Les books.}
\begin{enumerate}[resume]
\item Définir la \omterm{def:s2:dispersion-and-correlation:weighting}{pondération de la dispersion} et décrire les deux pondérations.
\item Donner les ratios de Sharpe et l'asymétrie des trois books.
\item Pourquoi le book sur les poids de l'indice perd-il en moyenne~?
\item Pourquoi le swap de corrélation est-il l'exposition la plus pure~?
\end{enumerate}

\medskip\noindent\textbf{Partie IV --- Le verdict.}
\begin{enumerate}[resume]
\item Énoncer le \emph{résultat nommé}~: le rendement du book de dispersion provenant de la prime de corrélation et sa perte lorsque la corrélation fait un saut.
\item Que s'est-il passé lors de chacune des trois crises~?
\item Pourquoi la troisième crise a-t-elle coûté moins cher~?
\item Comment dimensionner un book de dispersion~?
\item Comment faire un backtest honnête d'un tel book~?
\item Quelle fiche de stratégie est le plus exposé aux coûts des options sur titres individuels~?
\item Comment ce chapitre se rattache-t-il au chapitre~1~?
\item En une phrase~: que vend un trader de dispersion~?
\end{enumerate}
\end{problem}

\section{Questions d'entretien}

\begin{interviewq}[$\star$ \iqroles{chercheur}]\label{iq:s2:dispersion-and-correlation:1}
Qu'est-ce que la corrélation implicite, et comment la calcule-t-on~?
\end{interviewq}

\begin{interviewq}[$\star\star$ \iqroles{trader}]\label{iq:s2:dispersion-and-correlation:2}
Décrivez un trade de dispersion et les risques qu'il comporte.
\end{interviewq}

\begin{interviewq}[$\star\star$ \iqroles{chercheur}]\label{iq:s2:dispersion-and-correlation:3}
Pourquoi les options sur indice peuvent-elles être chères par rapport aux options sur titres individuels~?
\end{interviewq}

\begin{interviewq}[$\star\star$ \iqroles{risque}]\label{iq:s2:dispersion-and-correlation:4}
Comment soumettriez-vous un book de dispersion à un test de résistance~?
\end{interviewq}

\begin{interviewq}[$\star\star$ \iqroles{développeur}]\label{iq:s2:dispersion-and-correlation:5}
De quelles données un desk de dispersion a-t-il besoin chaque jour, et que peut-il arriver à ces données~?
\end{interviewq}

\begin{interviewq}[$\star\star\star$ \iqroles{chercheur}]\label{iq:s2:dispersion-and-correlation:6}
Pour $N$ membres à pondération égale, de volatilité $\sigma$ identique et de corrélation moyenne $\rho$, montrer que la variance de l'indice vaut $\sigma^2(\rho + (1 -
\rho)/N)$, puis en déduire la corrélation implicite à partir des volatilités de l'indice et des membres pour $N$ grand.
\end{interviewq}
